\section{Contexto del proyecto y restricciones}

El proyecto cuenta con un plazo estricto de \textbf{dos semanas} de desarrollo (iniciado el 22 de octubre de 2026), trabajado por tres estudiantes de la UNSAAC a tiempo parcial. Para cumplir con el tiempo, se seleccionaron herramientas estables y un flujo de trabajo ágil y ordenado.

\subsection{Tecnologías utilizadas}

\begin{itemize}[nosep]
    \item \textbf{Backend:} \textbf{Laravel 13} con \textbf{PHP 8.3}. Manejo modular de rutas, controladores, Form Requests tipados, políticas de autorización (RBAC) y el ORM Eloquent sobre MySQL.
    \item \textbf{Base de Datos:} \textbf{MySQL 8.0} (motor InnoDB, cotejamiento \texttt{utf8mb4\_unicode\_ci}). Integridad referencial con claves foráneas e índices únicos para control estricto de cupos y matrículas.
    \item \textbf{Frontend SPA:} \textbf{Inertia.js v3} con \textbf{Vue 3 (TypeScript)} y \textbf{Tailwind CSS}, ofreciendo una experiencia de aplicación de una sola página reactiva y fluida.
    \item \textbf{Servicios externos de identidad:} \textbf{API REST PeruDevs}, desacoplada mediante la interfaz \texttt{DniLookupService} para validación oficial de DNI ante el Registro Nacional de Identificación y Estado Civil (RENIEC).
    \item \textbf{Gestión del proyecto:} \textbf{Jira Software}. Tablero ágil para organizar el backlog de historias (\texttt{SIGC-XX}), limitar el trabajo en curso (WIP) y medir el avance del equipo.
    \item \textbf{Control de versiones y pruebas:} \textbf{GitHub} con flujo de ramas por desarrollador (\texttt{ramaFRANKI}), Pull Requests y suite de pruebas automatizadas con \textbf{Pest PHP} y \textbf{PHPUnit}.
\end{itemize}

\subsection{Stakeholders (Involucrados en el sistema)}

\begin{table}[H]
\centering
\small
\begin{tabularx}{\textwidth}{P{3.8cm} Y}
\toprule
\textbf{Involucrado} & \textbf{¿Qué hace o qué necesita en el sistema?} \\
\midrule
\textbf{Empresas / Instituciones} & Convocan a las capacitaciones y revisan los reportes finales de asistencia y aprobados. \\
\addlinespace
\textbf{Administrador} & Crea los cursos, asigna docentes, supervisa las vacantes y autoriza la emisión de certificados. \\
\addlinespace
\textbf{Docente / Instructor} & Dicta las clases, proyecta el código QR para la asistencia y sube las notas de los alumnos. \\
\addlinespace
\textbf{Participante / Alumno} & Se inscribe con su DNI, marca su asistencia escaneando el QR con su celular y descarga su certificado PDF. \\
\addlinespace
\textbf{Verificador externo} & Cualquier persona o empresa que escanea el código QR para comprobar si el certificado es original. \\
\bottomrule
\end{tabularx}
\caption{Principales involucrados en el sistema SIGC-CUSCO.}
\label{tab:stakeholders}
\end{table}
